\section{Morse attainment by handle cancellation}\label{app:handles}

We prove the existence assertion of Theorem~\ref{thm:smale} for a simply
connected closed smooth $D$-manifold $W$, with $D\geq9$.
The argument is the classical handle proof of Smale's theorem
\cite[Theorem~2.4]{Smale1963}.  Its purpose here is to keep the geometric
attainment of the integral torsion bound explicit.

We use the standard correspondence between Morse functions and ordered
handle decompositions, together with handle slides, insertion of
cancelling pairs, and handle cancellation.  A slide over a handle of the
same index changes the corresponding integral handle basis by an
elementary addition.  Two consecutive handles of indices $i$ and $i+1$
cancel when the attaching sphere of the latter meets the belt sphere of
the former transversely in one point.  Handle isotopies, slides, and
reorderings carry the subsequent attaching maps along.  Thus they preserve
the represented smooth manifold.  The Morse attachment correspondence is
reviewed in \cite{Milnor1963}.  We use the cancellation and basis-change
theorems in \cite[Theorems~5.4 and~7.6]{Milnor1965}.

\subsection{Removing handles of index one and coindex one}

Choose an ordered handle decomposition of $W$.  Connectedness allows us
to combine the $0$-handles along a tree of $1$-handles, leaving one
$0$-handle.  The dual operation leaves one $D$-handle.  We next remove
all $1$-handles, allowing the insertion of cancelling $2/3$ pairs.
This is the handle exchange underlying
\cite[Theorem~8.1]{Milnor1965}; we describe it to specify the role of
simple connectivity.

Let $P$ be the union of the $0$- and $1$-handles.  It retracts onto a
graph, and $\pi_1(P)$ is the free group generated by the $1$-handles.
The boundary $\partial P$ is connected, and its inclusion into $P$
induces an isomorphism on fundamental groups: viewed relative to
$\partial P$, the reversed handles have indices $D-1$ and $D$.
The attaching circles of the $2$-handles normally generate
$\pi_1(\partial P)$, because the higher-index handles do not affect
fundamental groups and $\pi_1(W)=0$.

Fix one free generator.  Insert a cancelling $2/3$ pair whose
$2$-handle initially attaches along a trivial circle in a small ball in
$\partial P$.  The new $2$-handle can be placed with the other
$2$-handles in the ordered decomposition.  Slide it over those handles
until its attaching circle represents the chosen generator.  To justify
this step, express the generator as a finite product of conjugates of
the old attaching words and their inverses.  A slide is a band sum
with an attaching circle, and choosing the band path and its orientation
realizes multiplication by the corresponding conjugate or inverse.
The needed paths can be chosen away from the other attaching circles by
general position; their tubular attaching regions can be made thinner.

The new circle is homotopic to a standard circle running once over the
chosen $1$-handle and avoiding the other $1$-handles.  Choose this standard
circle disjoint from the other $2$-handle attaching circles.  In the
$(D-1)$-dimensional boundary, general position turns the homotopy into an
isotopy through embedded circles avoiding those other circles.  Indeed,
both self-intersections in a one-parameter family and intersections with
the other circles have expected dimension at most $3-(D-1)<0$.
Transport the attaching framing along the isotopy.  The new attaching
circle now meets the chosen belt sphere once, so the $1$-handle and the
new $2$-handle cancel.  Repeating removes all $1$-handles.

Apply the same construction to the reversed decomposition to eliminate
the $(D-1)$-handles.  Its inserted $2/3$ pairs have original indices
$D-2,D-3$, so this operation does not reintroduce $1$-handles.  We
obtain an ordered decomposition with one handle in each of indices $0$
and $D$, none in indices $1$ and $D-1$, and all other handles in indices
$2,\ldots,D-2$.

\subsection{An algebraic unit whenever the count is excessive}

The simply connected manifold $W$ is orientable.  Fix an orientation,
orient the handles, and let $C_*$ be their integral chain complex.  The
coefficient of an $i$-handle in the boundary of an $(i+1)$-handle is the
algebraic intersection number of its belt sphere and the latter's
attaching sphere in the level between these indices.  Write
$Z_i=\ker\partial_i$ and $B_i=\operatorname{im}\partial_{i+1}$ as in
\eqref{eq:chain-rank}.  That identity gives
\begin{equation}\label{eq:handle-excess}
 \sum_i\rank C_i-\lambda(W)
 =2\sum_i\bigl(\rank B_i-t_i(W)\bigr).
\end{equation}
Each summand on the right is nonnegative.  If the total handle count
exceeds $\lambda(W)$, some $B_i$ therefore has rank greater than
$t_i(W)$.  Since the chain groups in degrees $1$ and $D-1$ vanish,
such an index satisfies $2\leq i\leq D-3$.

The subgroup $Z_i$ is a direct summand of $C_i$, because
$C_i/Z_i\cong B_{i-1}$ is free.  The nonzero Smith invariant factors of
$\partial_{i+1}$ are consequently those for the inclusion
$B_i\subset Z_i$.  Exactly $t_i(W)$ of them are nonunits: these are the
cyclic invariant factors of $\Tor H_i(W;\Z)$.  The strict rank inequality
thus gives at least one unit invariant factor.  Row and column operations
can place a $1$ in the boundary matrix.  They are realized geometrically
by handle slides, relabelings, and choices of orientation in indices $i$
and $i+1$.

The relevant levels are connected, so the slide bands can be chosen.
More precisely, the union $P_i$ of all handles of index at most $i$ is
simply connected, since it has no $1$-handles.  Its reversed handles
relative to the boundary $L=\partial P_i$ have indices at least
$D-i\geq3$.  They induce isomorphisms on components and fundamental
groups; hence $L$ is connected and simply connected.  The same reasoning
applies to the level before the $i$-handles.  Attaching spheres of one
index have codimension at least two in their attaching level, which
allows slide paths to avoid their disjoint attaching regions after those
regions have been made sufficiently thin.

We have produced an attaching sphere $A\cong S^i$ and a belt sphere
$B\cong S^{D-1-i}$ in $L$, with algebraic intersection number $1$.
To cancel their handles, we must make their geometric intersection a
single point.  This is the remaining geometric step.

\subsection{Converting the unit to a single intersection}

First suppose $i\geq3$ and $D-1-i\geq3$.  Put all attaching and belt
spheres at this level in transverse position.  A pair of opposite-sign
intersections of $A$ and $B$ determines a Whitney circle, using one arc
in each sphere and choosing the arcs to avoid the other intersection
points.  Since $L$ is simply connected, the circle bounds a Whitney
disk.  The high-dimensional Whitney trick
\cite[Theorem~6.6]{Milnor1965} removes the pair.  Its
dimensional hypotheses are precisely available here: both sphere
dimensions are at least three, and the interior of the two-dimensional
Whitney disk can avoid every attaching and belt sphere.  The Whitney
move can thus be performed without disturbing their other intersections.
Repeating leaves one intersection of $A$ and $B$.

The cases in which one sphere has dimension two require control of the
complement of the other sphere, as in the low-index refinement of the
Whitney theorem in \cite[Theorem~6.6]{Milnor1965}.  Here this control
can be seen directly.  By reversing the decomposition if necessary, it is enough to
treat $i=2$.  Here $L$ is the boundary after the $2$-handles have been
attached to the unique $0$-handle.  Let
\[
 B_1,\ldots,B_h\subset L
\]
be all their belt spheres, and put
$V=L\setminus\bigcup_{j=1}^h B_j$.
The plan is to move a disk in the chosen attaching sphere into $V$ by a
relative homotopy, then use general position to turn the resulting sphere
homotopy into an isotopy. We first determine the fundamental group of this
complement.

The manifold $L$ is obtained from $S^{D-1}$ by surgery on the attaching
circles.  Each surgery replaces
$S^1\times D^{D-2}$ by $D^2\times S^{D-3}$, whose belt sphere is
$\{0\}\times S^{D-3}$.  After removing this belt sphere, the replacement
piece is
\[
 (D^2\setminus\{0\})\times S^{D-3}.
\]
It retracts onto its attaching boundary
$S^1\times S^{D-3}$.  Performing these retractions simultaneously shows
that $V$ has the homotopy type of the exterior of the original attaching
circles in $S^{D-1}$.  That exterior is simply connected: a loop can be
filled in the sphere, and a filling disk can be put in general position
away from the circles.  Moving away from their cores and away from their
small tubular neighborhoods gives the same fundamental group.  Thus
\begin{equation}\label{eq:belt-complement-pi1}
 \pi_1(V)=\pi_1(L)=0.
\end{equation}

Let $A\cong S^2$ be the chosen attaching sphere of a $3$-handle, and
choose a pair of opposite-sign intersections with the chosen belt
sphere.  Join them by an arc in $A$ avoiding all other intersections
with all the $B_j$.  A small disk $\Delta\subset A$ around this arc
contains exactly that pair; its boundary and a collar outside it lie in
$V$.  The oriented normal bundles of the belt spheres have rank two.
Excision and the Thom isomorphism identify
\begin{equation}\label{eq:relative-intersection-group}
 H_2(L,V;\Z)\cong\bigoplus_{j=1}^h\Z,
\end{equation}
where the coordinates of $[\Delta]$ are its intersection numbers with
the $B_j$.  They all vanish, so $[\Delta]=0$.

The same conclusion holds in relative homotopy.  To see explicitly why
no assumption on $\pi_2(L)$ is needed, \eqref{eq:belt-complement-pi1}
and the long exact sequences give
\begin{align*}
 \pi_2(L,V)&\cong
 \operatorname{coker}\bigl(\pi_2(V)\longrightarrow\pi_2(L)\bigr),\\
 H_2(L,V;\Z)&\cong
 \operatorname{coker}\bigl(H_2(V;\Z)\longrightarrow H_2(L;\Z)\bigr).
\end{align*}
Absolute Hurewicz in degree two is an isomorphism for both simply
connected spaces.  It therefore induces an isomorphism between these
cokernels.  The disk $\Delta$, regarded as a map of pairs into $(L,V)$,
can consequently be homotoped into $V$, with its boundary moving in
$V$.  The homotopy extension property extends this boundary motion
through maps of the collar into $V$, fixing its outer edge.  Together
with the disk homotopy this gives a homotopy of $A$ stationary outside
$\Delta$ and the collar.  At its endpoint the
chosen pair has disappeared and every other intersection with the belt
spheres is unchanged.

We can realize this homotopy by an isotopy of the attaching sphere.
First smooth it relative to the initial sphere and the stationary
region.  General position for a one-parameter family of maps of a
$2$-sphere eliminates self-intersections and intersections with the
other attaching $2$-spheres: their expected dimensions are at most
$2+2+1-(D-1)<0$.  The immersion condition holds by the same parametric
general-position argument.  The level has dimension at least eight, so
these perturbations can be made relative to the prescribed region.
At the endpoint the altered disk and collar form a compact subset of
the open set $V$; a sufficiently small approximation preserves their
avoidance of the belt spheres.  Keep neighborhoods of the remaining
transverse intersections stationary.  Intermediate crossings with belt
spheres impose no restriction on this isotopy.  Isotopy extension
transports the attaching framing, and the other attaching tubes can be
shrunk to remain disjoint.  Repeating the operation leaves exactly one
intersection with the chosen belt sphere.

\subsection{Completing the cancellation}

In every case, the unit coefficient has now been represented by one
geometric intersection.  Order the chosen $i$-handle last among the
$i$-handles and the chosen $(i+1)$-handle first among the
$(i+1)$-handles.  They are consecutive and cancel.  Intersections with
other handles do not obstruct this cancellation; their attaching maps
are transported by the cancellation diffeomorphism.  The resulting
decomposition still has no handles of indices $1$ or $D-1$.

Every cancellation decreases the total number of handles by two.
Whenever that count exceeds $\lambda(W)$, Equation
\eqref{eq:handle-excess} produces another cancellable pair by the
argument above.  The process therefore terminates with exactly
$\lambda(W)$ handles.  At this stage every nonnegative summand in
\eqref{eq:handle-excess} vanishes, so $\rank B_i=t_i(W)$ for every $i$.
Equation~\eqref{eq:chain-rank} then gives precisely
\[
 \rank C_i=r_i(W)+t_i(W)+t_{i-1}(W).
\]
The Morse function associated with the final handle decomposition has
these critical-point counts.  This proves the existence assertion in
Theorem~\ref{thm:smale}.
